%  Bachelorarbeit                                           %%
%% TH Köln -Campus Gummersbach, Fak. 10)                    %%
%% 2021                                                     %%
\documentclass[a4paper,12pt,oneside]{article}
% Optionen:
% - a4paper => DIN A4-Format
% - 12pt    => Schriftgröße (weitere  
%              grundlegende Fontgrößen: 10pt, 11pt)
% - oneside => Einseitiger Druck

%% Verwendete Pakete:
\usepackage[ngerman]{babel} % für die deutsche Sprache
\usepackage{caption} % Für schönere Bildunterschriften
\usepackage[T1]{fontenc} % Schriftkodierung (Für Sonderzeichen u.a.)
\usepackage[utf8]{inputenc} % Für die direkte Eingabe von Umlauten im Editor u.a.
\usepackage{fancyhdr} % Für Kopf- und Fußzeilen
\usepackage{lscape} % Für Querformat

%% Show code in a better way
\usepackage{listings}
\usepackage{color}

\definecolor{dkgreen}{rgb}{0,0.6,0}
\definecolor{gray}{rgb}{0.5,0.5,0.5}
\definecolor{mauve}{rgb}{0.58,0,0.82}

\lstset{frame=tb,
  language=Java,
  aboveskip=3mm,
  belowskip=3mm,
  showstringspaces=false,
  columns=flexible,
  basicstyle={\small\ttfamily},
  numbers=none,
  numberstyle=\tiny\color{gray},
  keywordstyle=\color{blue},
  commentstyle=\color{dkgreen},
  stringstyle=\color{mauve},
  breaklines=true,
  breakatwhitespace=true,
  tabsize=3
}
%% Schriften (Beispiele)
%% Weitere LaTeX-Schriften im "LaTeX Font Catalogue"
%% unter: http://www.tug.dk/FontCatalogue/.
%% ACHTUNG: Ggf. müssen Schriften noch installiert 
%% werden!

% Serifen-Schriften:
\usepackage{lmodern} % Schriftart "Latin Modern"
%\usepackage{garamond} % Schriftart "Garamond"

%Sans Serif-Schriften:
%\usepackage[scaled]{uarial}
%\usepackage[scaled]{helvet}
%%--------------
\usepackage[normalem]{ulem} % Für das Unterstreichen von Text z.B. mit \uline{}
\usepackage[left=3cm,right=2cm,top=1.5cm,bottom=1cm,
textheight=245mm,textwidth=160mm,includeheadfoot,headsep=1cm,
footskip=1cm,headheight=14.599pt]{geometry} % Einrichtung der Seite 

\usepackage{graphicx} % Zum Laden von Graphiken
\graphicspath{ {./sources/} }

% INFO: Graphiken einbinden
%
% \includegraphics[scale=1.00]{dateiname}
%
% => Ausgabeformat: PDF-Dokument:
%    Es können die folgenden (Graphik-)formate eingebunden
%    werden: .jpg, .png, .pdf, .mps
% 
% => Ausgabeformat: DVI/PS:
%    Folgende (Graphik-)formate werden unterstützt:
%    .eps, .ps, .bmp, .pict, .pntg
\usepackage{epstopdf}

% Pakete für Tabellen
\usepackage{tabularx} % Einfache Tabellen
\usepackage{longtable} % Tabellen als Gleitobjekte (für die Aufteilung bei langen 
 %Tabellen über mehrere Seiten)
\usepackage{multirow} % Für das Verbinden von Zeilen innerhalb einer Tabelle mit
 % \multirow{anzahl}{*}{Text}

% (Zusatz-)Pakete für Formeln
\usepackage{amsmath}
\usepackage{amsthm}
\usepackage{amsfonts}

\usepackage{setspace} % Paket zum Setzen des Zeilenabstandes
% INFO: Zeilenabstand setzen:
%
% Befehle:
% - \singlespacing  => 1-zeilig (Standard)
% - \onehalfspacing => 1,5-zeilig
% - \doublespacing  => 2-zeilig 
\onehalfspacing % Zeilenabstand auf 1,5-zeilig setzen

% Farbboxen (für die Merkkästen in dieser Vorlage):
\usepackage{tcolorbox}
\tcbset{colback=white,colframe=orange,
        fonttitle=\bfseries}

\usepackage[colorlinks,pdfpagelabels,pdfstartview=FitH,
bookmarksopen=true,bookmarksnumbered=true,linkcolor=black,
plainpages=false,hypertexnames=false,citecolor=black]{hyperref} % Für Verlinkungen
% INFO: Verlinkungen mit dem hyperref-Paket:
%
% Die Angabe von URLs mit dem Befehl \url{} erlaubt einen
% gesonderten Umgang mit Weblinks. Denn die Links werden verlinkt.
% Auch erfolgt automatisch am Zeilenende ein Umbruch des Links.
% Es ist auch nicht erforderlich, Sonderzeichen in der URL manuell zu 
% entschärfen.
%
% TIPP: Sollte ein Umbuch bei einem Link nicht automatisch erfolgen, so kann
% das daran liegen, dass ein/mehrere Zeichen zusätzlich angegeben werden müssen,
% an dem der Link umbrochen werden kann.
% Dies kann mit folgendem Befehl erfolgen (Beispiel):
% \renewcommand*\UrlBreaks{\do-\do_}

% Das Paket "biblatex" für autom. 
% Literaturverzeichnisse:
%\usepackage{csquotes} % Für sprachangepasste Anführungszeichen
%\usepackage[backend=bibtex,style=alphabetic]  
%           {biblatex}
%\addbibresource{bib/literatur.bib}           

%%%%%%%%%%%%%%%%%%%%%%%%%%%%%%%%%%%%%%%%%%%%%
%% DOKUMENT                                %%
%%%%%%%%%%%%%%%%%%%%%%%%%%%%%%%%%%%%%%%%%%%%%

\begin{document}
% Unbeschriftetes Vorblatt (Leere Seite)
\pagestyle{empty} % Seite ohne Kopf- und Fußzeilen
\newpage % Neue Seite
\input{leereSeite} % Ausgelagerte LaTeX-Datei (hier: leereSeite.tex) einbinden

\newpage

% Deckblatt
\pagestyle{empty}
\begin{titlepage}
  \includegraphics[scale=0.20]{sources/TH_Koeln_Logo}\\
  \begin{center}
    \Large
    Technische Hochschule Köln\\
    Fakultät für Informatik und Ingenieurwissenschaften\\
    \hrule\par\rule{0pt}{2cm} % Horizontale Trennlinie  mit 2 cm Abtand nach unten erzeugen
    \LARGE
    \textsc{BACHELORARBEIT}\\
    \vspace{1cm} % Vertikaler Abstand von 1cm erzeugen
    \huge
    Kostenoptimierung für AWS Cloud-Diensten \\
    \Large
    Ein wirtschaftlicher Ansatz für die Cloud\\
    \vspace{1cm}
    \large
    Vorgelegt an der TH Köln Campus Gummersbach\\
    im Studiengang Wirtshaftsinformatik\\
    \vspace{1.0cm}
    ausgearbeitet von:\\
    \textsc{Carlo Menjivar} 11117929\\
    \vspace{1cm}
    \begin{tabular}{ll} % Einfache Tabelle ohne Rahmen, mit 2 Spalten erzeugen
      \textbf{Erster Prüfer:}  & Prof. Dr. Roman Majewski \\
      \textbf{Zweiter Prüfer:} & <Name des 2. Prüfers>    \\
    \end{tabular}
    \vspace{1cm}
    \\Gummersbach, im F<Monat der Abgabe>
  \end{center}
\end{titlepage}

\newpage

% Abstract (ACHTUNG: Abweichung zur Reihenfolge im Merkblatt!)
\begin{abstract}
  Platz für das deutsche Abstract...
\end{abstract}

\renewcommand{\abstractname}{Abstract}
\begin{abstract}
  Platz für das englische Abstract...
\end{abstract}
%<MERKKASTEN> (für die eigene Verwendung bitte entfernen
\vspace{1cm}
\begin{tcolorbox}[title={Das Abstract}]
  Bei einem Abstract handelt es sich um eine Art \textit{Zusammenfassung} Ihrer Arbeit. Diese kann in deutscher und/oder englischer Sprache verfasst werden. Mithilfe des Abstracts kann der Leser sich zügig orientieren, in wie fern Ihre Arbeit für ihn Relevanz besitzt.\\                                                                      Sprechen Sie unbedingt mit Ihrer Betreuerin/Ihrem Betreuer, ob Sie für Ihre Arbeit ein Abstract benötigen.\\
  Ein Abstract beinhaltet folgende Aspekte \footnote{ Vgl. \cite{SW11}, S. 249}:
  \begin{itemize}
    \item Ziel der Arbeit
    \item Fragestellung der Arbeit
    \item Herangezogener, theoretischer Ansatz ("Quellen")
    \item \textit{Optional:} Methodik
  \end{itemize}
\end{tcolorbox}
%</MERKKASTEN>

%<MERKKASTEN> (für die eigene Verwendung bitte entfernen
\vspace{1cm}
\begin{tcolorbox}[title={Hinweise zu dieser Dokumentvorlage}]
  \begin{itemize}
    \item Es handelt sich hierbei um eine Beispiel-Vorlage für wissenschaftliche Ausarbeitungen.
          Über die konkrete, formale Ausgestaltung Ihrer wissenschaftlichen Arbeit sprechen Sie unbedingt mit Ihre/m Betreuer/in.
    \item Unabhängig, ob Sie beispielsweise eine Bachelor-, Master- oder Hausarbeit schreiben müssen. Diese Vorlage kann als eine gute Basis für Ihre Arbeit dienen. Passen Sie einfach die Vorlage Ihren Anforderungen entsprechend an.
  \end{itemize}
\end{tcolorbox}
%</MERKKASTEN>  

\newpage

% Inhaltsverzeichnis
\tableofcontents

\newpage
\pagestyle{fancy} % Kopf- und Fußzeilen aktivieren (=> Paket "fancyhdr")

% Abbildungsverzeichnis  
% INFO: Abbildung einbinden (Beispiel):
%  \begin{figure}[h!]
%    \centering
%    \includegraphics[scale=1.00]{Pfad zum Bild}\\
%    \caption{Bildunterschrift} 
%    \label{Marke zum Referenzieren auf die Abbildung}
%  \end{figure}
\section*{Abbildungsverzeichnis}
\addcontentsline{toc}{section}{Abbildungsverzeichnis} % Manuellen Eintrag im Inhaltsverzeichnis erzeugen
\renewcommand{\listfigurename}{} % Name des Abbildungsverzeichnisses ändern
\thispagestyle{empty}
\listoffigures

\newpage
% INFO: Querverweise auf Gliederungselemente, Abbildungen 
%       & Tabellen setzen:
%
% Voraussetzung: Gesetzte Referenzmarke mit dem Befehl: \label{marke}
% 
% Referenzierung erfolgt dann mittels dem Befehl:
% \ref{marke}

\section{Problemstellung}\label{kap_problemstellung}
% Was ist das Problem
\paragraph{}
Unsere Aufgabe ist es, eine Webplattform für Spieler des Videospiels "League of Legends" zu schaffen, auf der sich Spieler kennenlernen und zu einer gemeinsamen Partie verabreden können. Um eine möglichst gute Plattform zu bieten, sind verschiedene Schritte zu erledigen.

\subsection{Aufmerksamkeit/Marketing}
\paragraph{}
Im Rahmen der Projektarbeit steht Marketing nicht im Fokus. Wir beschreiben daher nur kurz, worauf zu achten ist, wenn das Produkt nach dem Praxisprojekt auf den Markt kommen soll.

\paragraph{}
Wie bei sozialen Plattformen und sozialen Netzwerken üblich, erhöht sich der Nutzen durch positive Netzwerkeffekte mit steigender Nutzeranzahl. Ein Ziel ist es daher, eine möglichst große Nutzerbasis aufzubauen, die unsere Webseite möglichst oft verwendet.
Zuerst müssen wir Mögichkeiten finden, damit Personen von unserer Webseite überhaupt erfahren. Auf technischer Seite ist dies mit Suchmaschinenoptimierung möglich, Marketingtechnisch haben wir verschiedene Möglichkeiten der Werbung, unter anderem Bannerwerbung, Influencermarketing und "Kunden werben Kunden".

Sobald Personen unsere Webseite besuchen, müssen diese überzeugt werden, dass unser Produkt gut ist und dass sie ein Konto erstellen sollten. Es sollte daher auf einer ansprechenden Startseite das Produkt vorgestellt werden, auf dem erklärt wird, was das Ziel unserer Plattform ist und wie wir dieses Ziel erreichen wollen.

Je besser unser Produkt ist, desto wahrscheinlicher ist es, dass ein Nutzer durch intrinsische Motivation unsere Webseite verwendet. Es ist daher wichtig, ein Produkt mit möglichst hoher Qualität zu erstellen. Ein Nutzer, der von unserer Webseite begeistert ist, wird wahrscheinlicher die Webseite öfter verwenden und Freunden von der Webseite erzählen, was wiederum neue Nutzer generieren kann.

\subsection{Das Nutzerkonto}
\paragraph{}
Um die Webseite vollständig nutzen zu können, sollen Besucher der Seite ein Konto erstellen. Dazu sollen sie eine E-Mail angeben sowie ein Passwort und einen freigewählten Nutzernamen aussuchen. Auch sollen sie nach Kontoerstellung Nutzerdaten wie ein Profilbild, das Alter und einen Freitext angeben können. Die angegeben Daten sollen dem Nutzer dabei helfen, Kontakte auf der Plattform zu finden und können als Gesprächsthema dienen.

Auch ist geplant, dass das Nutzerkonto mit dem Riot-Konto verbunden werden kann. Dies soll es ermöglichen, dass Spielstatistiken, wie zum Beispiel die meist gespielten Champions, die Lieblingsposition und die Elo, angezeigt werden. Ein Nutzer, der auf Kontaktsuche ist, erhält damit mehr Informationen über den Spieler und kann somit besser abschätzen, ob er Kontakt aufnehmen will. 

\subsection{Kontaktsuche}
\paragraph{}
Um mit anderen Nutzern in Kontakt treten zu können, soll es möglich sein, andere, zufällige Nutzerprofile anzuzeigen. Profile sollen einzeln hintereinander angezeigt werden; nachdem sich der Nutzer entschieden hat, ob er dem angezeigten Nutzer einen Like [GLOSSAR!] vergeben will, wird das nächste Profil angezeigt.

Sobald zwei Nutzer gegenseitig einen Like vergeben haben, sollen diese befreundet sein. Befreundete Nutzer sollen untereinander Nachrichten austauschen können.

\subsection{Filter}
\paragraph{}
Um die Qualität der angezeigten Profile zu erhöhen und eine möglichst gute Nutzererfahrung zu gewährleisten, soll es dem Nutzer möglich sein, seine Suche einzugrenzen. Durch Filter werden ihm nur die Nutzer angezeigt, die für ihn relevant sind.

\subsection{Regeln und Richtlinien}
\paragraph{}
Wir möchten mit unserer Plattform einen sicheren Ort bieten, auf dem die Spieler unabhängig ihrer Eigenschaften respektiert werden. Wir sind der Meinung, dass die meisten Nutzer in der Lage sind, respektvoll miteinander umzugehen.
Dies können wir jedoch nicht gewährleisten.

Wir wollen uns daher das Recht vorbehalten, Nutzern, die gegen unsere Regeln und Richtlinien verstoßen, den Zugang zu unserer Plattform zu verwehren. Sollte sich ein Nutzer von einem anderen Nutzer beleidigt fühlen oder andersweitig der Meinung sein, dass dieser Nutzer gegen die Regeln verstößt, soll dieser Regelbrecher gemeldet werden können. Wir als Betreiber der Plattform sollen in der Lage sein, die Meldungen zu verarbeiten und angemessen mit den Regelbrechern zu verfahren.

Unabhängig davon soll es den Nutzern möglich sein, andere Nutzer zu blockieren. Blockierte Nutzer sollen dem aktiven Nutzer nicht mehr in der Spielersuche angezeigt werden, nicht mehr Nachrichten verschicken und aktive Chats sollen aufgelöst werden. Einen anderen Spieler zu blockieren heißt nicht zwangsläufig, dass dieser gegen Regeln oder Richtlinien verstößt, es kann auch sein, dass der Nutzer einfach den Kontakt abbrechen will. Sollte ein Nutzer jedoch gemeldet werden, soll er automatisch blockiert werden.

\subsection{Internationalisierung}
Die Webseite soll in Englisch angeboten werden, da dies eine der weit verbreitetsten Sprachen der Welt ist und damit viele Nutzer erreicht.

\subsection{Einstellungen}
Um den verschiedenen Vorlieben unterschiedlichster Nutzer gerecht zu werden, sollen diese in der Lage sein, verschiedene Optionen in den Einstellungen zu verwalten. Unter anderem sollen Nutzer in der Lage sein, auszusuchen, wie und ob sie über neue Nachrichten, Freundschaften und Neuigkeiten informiert werden.

\subsection{Nachrichten}
Sobald zwei Nutzer befreundet sind, können sie sich gegenseitig Nachrichten schreiben. Dazu soll auf einer dafür eingerichteten Seite in ein Eingabefeld der zu versendende Text eingegeben werden. Sobald die Nachricht abgeschickt wird, soll der andere Nutzer möglichst Zeitnah von dieser erfahren. Dies soll den glatten Ablauf von Chats ermöglichen.

Zeichen sollen nach dem UTF-8 Zeichensatz erlaubt sein.

\subsection{Sicherheit}
Die Sicherheit der Nutzerdaten ist von hoher Priorität. Passwörter müssen besonders geschützt und nach derzeitigem kryptologischen Stand der Technik gesichert werden, um den Zugriff von unautorisierten Angreifern zu unterbinden. Identifizierende Daten wie E-Mail-Adressen dürfen nicht öffentlich einsehbar sein.

Sollte in einem späteren Schritt die Plattform veröffentlicht werden, ist auf DSGVO-Konformität zu achten.

\subsection{Dokumentation}
"Technische Schulden" (en. "technical dept", schlechte Umsetzung von Software, die einen erheblichen Mehraufwand in der Zukunft bedeutet) sorgen in vielen Fällen für Probleme im Laufe des Produktlebenszyklus. Um die Wahrscheinlichkeit zu verringern, dass sich technische Schulden anhäufen und, soll eine gut leserliche Dokumentation zum Quellcode geschrieben werden, welche es allen Gruppenmitgliedern ermöglichen soll, den überblick zu halten, welcher Codeschnipsel welches Problem löst. Desweiteren sollen umfangreiche Tests durchgeführt werden, welche die Qualität des Codes gewährleisten sollen.

Um die Kommunikation zwischen Backend und Frontend zu gewährleisten, sollen Schnittstellen entsprechend dokumentiert werden. Es sollte einem Frontend-Entwickler möglich sein, nur mit der Dokumentation auf die Schnittstelle zuzugreifen und die Informationen zu erhalten, die er benötigt, um entsprechende Daten im Frontend anzeigen zu können.

\subsection{Konto löschen}
Es wird Benutzer geben, die aus verschiedensten Gründen unsere Webseite nicht weiter verwenden wollen und ihr Konto löschen wollen. Unsere Aufgabe ist es, eine sichere Löschung des Benutzerkontos zu gewährleisten und persönliche Daten nach Datenschutzvorschriften aus der Datenbank zu entfernen.

Optional können wir zudem den Nutzer darum bitten, ein Formular auszufüllen, in dem dieser uns Rückmeldung geben kann, aus welchen Gründen er die Webseite nicht weiter verwenden will.


\newpage
\section{Einleitung}\label{kap_einleitung}
%Motivation und Ziele als subsecion?
%WEITERE QUELLE Für mOTIVATION
%https://www.gartner.com/smarterwithgartner/4-trends-impacting-cloud-adoption-in-2020
%\subsection*{Einführung in das Thema (Motivation, zentrale Begriffe etc.)}
\subsection{Motivation}
%\addcontentsline{toc}{subsection}{Motivation}
%WARUM HABE ICH MICH FÜR AWS ENTSCHIEDEN?
%Wie viele Firmen wechseln von On-Premise zu Cloud in DE jährlich?
[Persönliches Statement... wieso schreib ich das hier eigentlich ;)] Die zunehmende Digitalisierung von Geschäftsmodellen, die auch durch die Corona-Pandemie vorangetrieben wird, lässt Cloud-basierte Applikationen an Bedeutung gewinnen.\footnote{Es sei an diser Stelle darauf hingewiesen, dass in diesem Kontext Ahrens die Bedeutung Cloud-basierter Anwendungen im Bereich von deutschen Handelsunternehmen untersuchte (Vgl. Ahrens 2021)\cite{STA3}, sowie das \textit{ifo Institut} anschaulich die strukturellen Veränderungen von der Corona-Pandemie auf den Arbeitsalltag in Deutschland nachzeichnete (Vgl. ifo Institut 2020)\cite{STA2}.} Als direkte Folge davon ist die Nachfrage nach Server- und Speicherkapazität gestiegen.
%Dies wird auch von 48\% der befragten Handelsunternehmen im Jahr 2021 bestätigt.
Die Relevanz von \textit{Amazon Web Services}, kurz AWS, in dem Bereich der Cloud-Computing ergibt sich aus einer vor kurzem veröffentlichte Studie von Raj Bala et al.. Diese wies eindrücklich daraufhin, dass AWS der aktuell weltweit führende Cloud-Anbieter anhand ihrer Klassifikation (\textit{Magic Quadrant}\footnote{ Laut Gartner stellt der Magic Quadrant eine zweidimensionale Matrix mit vier Quadranten dar. Jeder Quadrant steht für einen Unternehmenstypus im Markt. Im Uhrzeigersinn von links unten beginnend sind dies: \textit{Nichenanbieter, Herausforderer, Marktführer }und \textit{Visionäre}}) für Cloud-Infrastruktur und Plattform-Services sei (Bala et al, 2021, o.S.,\cite{G01}).
AWS erscheint nicht nur aus diesem Grund als Fallbeispiel für diese Arbeit passend, weitere bedeutsame Faktoren sind seine frühe Präsenz (2006) als Cloudanbieter und seines großen Angebotes an Cloud-Diensten, welche für zahlreiche Anwendungsfälle geeignet sind.\footnote{ Die aktuellen Marktführer im Bereich der \textit{Cloud-Computing} weltweit sind AWS, Google, Telekom und Microsoft (Vgl. Synergy Reseach Group 2019, o.S.\cite{STA6})}
\\\\

\begin{comment} GELÖSCHT, WEIL DIESE EINE BEHAUPTUNG IST (25.10.2021)
    \\\\
    Für viele Unternehmen ist eine große Herausforderung, die Kosten von Cloud-Diensten übersichtlich zu halten und Optimierungsmöglichkeit leicht zu erkennen. Zusätzlich besteht die Gefahr, unangenehme Überraschungen in einer Rechnung zu bekommen, weil keine Grenze für den Konsum von Cloud-Diensten festgelegt wurde. 
    \end{comment}
\subsection{Problemstellung}
%\addcontentsline{toc}{subsection}{Problemstellung}
%Wenn ein Hotel die Vorteile von dem Cloud-Computing hätte, dann könnte dieses folgendermaßen funktionieren:
\begin{comment}
\\\\
”Heute hatten wir 17 Gäste für unsere derzeit 20 Zimmer. Für die kommende Messe am Wochenende sind wir bereit 500 Gäste zu empfangen. Nach der Messe werden wir mit unseren üblichen 20 Zimmern wie immer gut arbeiten können.”
Normalerweise bräuchte man eine große Investition zu machen, um solche kurzfristige Nachfrage zu erfüllen. Vergleichbar ist es bei traditionellen IT-Infrastrukturen, mehr Kapazitätsbedarf, würde die Anschaffung von einer neuen Hardware bedeuten.
\\\\
\end{comment}
%Die Verwendung von Cloud-Diensten bringt viele Vorteile mit sich. Zum Beispiel kurzfristige Erhöhung oder Verringerung der Speicher- und Rechenkapazität, sowie Zugriff auf unterschiedliche Speicherarten, die genau an individuelle Anwendungsfälle angepasst sind. All diese Lösungen sind in wenigen Minuten einsatzfertig. 
%\\\\
%Viele Unternehmen befürchten jedoch, dass der Wechsel von On-Premise zu On-Demand zu hohen Kosten führen könnte.
%In einer Umfrage haben circa 50\% der Unternehmen die Verwaltung der Kosten für den Betrieb von Cloud-Workloads als großes Hindernis genannt. Mehr als die Hälfte der Befragten haben geäußert, dass sie Schwierigkeiten haben, alle Kosten für Cloud-Workloads zu rechtfertigen.
Adam Stern wies in dem \textit{Forbes}-Magazin daraufhin, dass ungefähr die Hälfte der US-amerikanischen Unternehmen Schwierigkeiten hätten ihre Kosten zu begründen (Stern 2018, o.S.). 
\begin{quote}
    „In its Stratecast Predictions 2018, Frost \& Sullivan noted that 53\% of IT leaders surveyed cited “managing costs to run cloud workloads” as a huge obstacle, and over 50\% have difficulty justifying the expenses of some public cloud workloads.“  
    \footnote{Stern, Adam, The Truth About Cloud Pricing.\cite{SP1}}
\end{quote}
Darüber hinaus weist Tobias Regenfuß und Jochen Malinowski
von Accenture GmbH in einer Untersuchung, dass es den Unternehmen an fachlichem Know-How in Cloud-Computing mangelte. Diese stelle eine der größten Hindernisse dar, um einen Wechsel von On-Premise- zu Cloud-basierten Systemen gewährleisten zu können\footnote{Regenfuß und MalinowskiStern, Hohe Erwartungen an die Cloud: Hürden meistern, Mehrwert maximieren. 2020 o.S.(Webversion) oder S.11 in der PDF-Version auf Englisch\cite{ACC1}}.
\\\\
Kostenoptimierung für Cloud-Dienste ist ein wichtiger Punkt, da man ohne Optimierungsmaßnahmen mit höheren Kosten rechnen müsse als bei On-Premise Systemen(Anders Lisdorf\footnote{Anders Lisdorf. Cloud Computing Basics: a Non.-Technical Introduction. S.152. \cite{CCB}}).
\\
([Rev] SOLLTE DIE UNTERE DIREKTE ZITAT WEG)
\begin{quote}
    ”Indeed, if you run the cloud the same way you run your on-premise data center, you are almost certain to incur higher expenses. It is necessary to use the following key cloud cost optimization techniques in order to successfully save money on the cloud.”
    \footnote{Anders Lisdorf. Cloud Computing Basics: a Non.-Technical Introduction. S.152. \cite{CCB}}
\end{quote}
\begin{flushleft}
Diese Bachelorarbeit beschäftigt sich mit ebendieser Problematik, um herauszufinden, wie Unternehmen mit den passenden Werkzeugen die Kosten ihrer Cloud-Dienste überwachen und im Blick behalten können. %Zum Beispiel können frühzeitige Benachrichtigungen alarmieren, wenn Cloud-Dienste mehr Kosten verursachen als geplant.
\\
Außerdem sollte untersucht werden, wie mit der richtigen Auswahl an Diensten Kosten optimiert werden. 
%In dieser Arbeit wird versucht zu beantworten, wie Kosten bei Cloud-Diensten überwacht werden können. Auf Grundlage dieser Information werden Optimierungsmöglichkeiten untersucht. 
Es wird untersucht, welche Maßnahmen nötig sind, um unerwartet hohe Kosten bei Cloud-Diensten zu vermeiden. Darüber hinaus werden Empfehlungen von Cloud-Experten berücksichtigt, um Kosten von Cloud-Diensten zu minimieren. Diese Arbeit untersucht speziell die Kostenoptimierung  von S3-Speichereinheiten und EC2-Server-Instanzen mithilfe von folgenden Überwachungswerkzeuge: Cost-Explorer, CloudWatch und Trusted Advisor.
\end{flushleft}


\subsection{Zielsetzung}
%\addcontentsline{toc}{subsection}{Zielsetzung}
Die vorliegende Arbeit betrachtet die von AWS angebotenen Überwachungswerkzeuge, um ein tiefergehendes Verständnis der Entstehung von Kosten durch die Nutzung von Cloud-Diensten zu gewährleisten. Mit den von AWS zur Verfügung gestellten Maßnahmen sollen die Nutzung und damit die Kosten von Cloud-Diensten reduziert werden.
%\textbf{Daraus ergeben sich für die Arbeit die folgenden Ziele:}%\\ 
%\%begin{itemize}
 %   \item
 %       Als Erstes wird gezeigt, wie mithilfe von bestehenden %Werkzeugen  die Kosten von Cloud-Diensten überwacht %werden können.
  %  \item
 %       Als Nächstes wird anhand von Empfehlungen von Cloud-Experten identifiziert, welche Optimierungsmöglichkeiten bestehen.\\
%\end{itemize}

%Die vorgestellten Werkzeuge werden auf eine Testumgebung eingesetzt und deren Auswirkungen im Bezug auf die Kosten %bewertet.\\
\begin{comment}
\subsection*{Einschränkungen}
\addcontentsline{toc}{subsection}{Einschränkungen}

Der Schwerpunkt dieser Arbeit liegt auf EC2-Instanzen, da diese in der Regel den größten Anteil an der Rechnung ausmachen.
An zweiter Stelle stehen S3-Speichereinheiten, weil sie einen erheblichen Teil der Kosten darstellen.
%STATISTEN DIE DAS BELEGEN?

%Nach Angaben von Amazon Web Services ist es möglich bis zu 90% für EC2 zu sparen, wenn EC2 Spot-Instanzen benutzt werden. 
%Eine Preisreduzierung für Speichereinheiten ist möglich, wenn die richtige Speicherart ausgewählt wird. 
{\cite{AMZ08,AMZ09}} 
\\
Diese Arbeit legt den Fokus auf die Optimierung der oben genannten Dienste.
Als Überwachungswerkzeuge für die Kosten werden die AWS CloudWatch, der AWS Cost-Explorer und der AWS Trusted Advisor untersucht. 
\end{comment}
\subsection{Struktur der Arbeit}
%\addcontentsline{toc}{subsection}{Struktur der Arbeit}

%Schlüsselbegriffe
%Zunächst wird es eine kurze Einführung in die relevanten Begriffe geben, die für die zu untersuchenden AWS-Cloud-Services wichtig sind.

Diese Bachelorarbeit ist in folgende Kapitel unterteilt:\\\\
%2 Die gängigsten Cloud-Dienste, bei deren Geld verschwendet wird.
\textbf{Kapitel~\ref{kap_grundlagen}} 
befasst sich mit dem Begriff Cloud-Economy und erläutert das Potenzial der Cloud-Diensten im wirtschaftlichen Sinne. Die Cloud-Dienste EC2-Instanzen und S3 Speichereinheiten werden ebenfalls kurz erklärt. %Diese dienen als Grundlage für diese Arbeit. 
\\\\
%4 Überwachung von Kosten
\textbf{Kapitel~\ref{kap_zahlungsmodelle}} 
zeigt die verschiedenen Zahlungsmodelle für EC2-Instanzen. Es werden Kriterien vorgestellt, die helfen sollen, sich für das richtige Zahlungsmodell bei verschiedenen Szenarien zu entscheiden. 
\\\\
%5 AWS Cost Explorer und AWS-Kosten- und Nutzungsbericht
\textbf{In Kapitel~\ref{kap_kostenüberwachung }} werden die Werkzeuge eingeführt, die zur Überwachung der Kosten von Cloud-Diensten eingesetzt werden.
\\\\
%4 Methoden zur Kostenbremse
%4.1 S3 Intelligent-Tiering
%4.2 Instance-Scheduler für EC2 und AWS Reserved-Instance
\textbf{Kapitel~\ref{kap_Optimierung}} befasst sich mit Optimierungsmaßnahmen %insbesondere 
für EC2-Instanzen und S3 Speichereinheiten.

%Testumgebung
%Schließlich werden anhand eines Fallbeispiels in \textbf{Kapitel 5?}, die oben genannten Werkzeuge und Techniken in einer %kostenlosen Testumgebung getestet. Um die Zuverlässigkeit der Ergebnisse zu gewährleisten, wird alles gemacht, um die %Vorher- und Nachher-Szenarien vergleichbar zu machen. Dabei werden die Anzahl der Instanzen und deren Auslastung sowie %die Daten auf den Speichereinheiten berücksichtigt.
 


\vspace{1cm}
\begin{tcolorbox}[title={Die Einleitung(DELETE AFTER CHECK WE COMPLETE ALL POINTS)}]
  Die Einleitung umfasst folgende Elemente\footnote{Vgl. u.a. \cite{BBoJ}, S. 5-6}:
  \begin{itemize}
    \item Einführung in das Thema (Motivation, zentrale Begriffe etc.)
    \item Hinführung zu den Ergebnissen
    \item Ggf. Angabe des Schwerpunktes
    \item Ggf. Einschränkungen darlegen
    \item Problemstellung
    \item Zielstellung der Arbeit
    \item Fragestellung der Arbeit
    \item Übersicht über die Kapitel geben:
          \begin{quotation}
            Eine Einleitung muss auch durch die Arbeit führen. Sie muss dem Leser helfen, sich in der Arbeit und ihrer Struktur zu Recht zu finden. Für jedes Kapitel sollte eine ganz kurze Inhaltsangabe gemacht werden und ggf. motiviert werden, warum es geschrieben wurde. Oft denkt sich ein Autor etwas bei der Struktur seiner Arbeit, auch solche Beweggründe sind dem Leser zu erklären\footnote{\cite{BBoJ}, S. 6}:.
          \end{quotation}
  \end{itemize}
\end{tcolorbox}

\newpage
\section{Grundlagen}\label{kap_grundlagen}
TEXT FOLGT...

\subsection{Unterabschnitt von Grundlagen}\label{subsec_UabsGrundl}
%\input{}
TEXT FOLGT...

%<MERKKASTEN> (für die eigene Verwendung bitte entfernen
\vspace{1cm}
\begin{tcolorbox}[title={Das Kapitel/der Abschnitt}]
  Hierbei handelt es sich um ein Beispiel-Kapitel. Es ist zu empfehlen, dass Sie Kapitel und auch Abschnitte immer mit einer kurzen Einleitung beginnen. In dieser beschreiben Sie kurz, was den Leser in diesem Kapitel/Abschnitt erwartet. Bei einem Kapitel mit Abschnitten nehmen Sie auch inhaltlichen Bezug auf die enthaltenen Abschnitte (inklusive Referenzierung auf die Abschnittsnummerierung).
\end{tcolorbox}
%</MERKKASTEN>  

\vspace{1cm}
\begin{tcolorbox}[title={Abbildungen, Tabellen \& Co.}]
  Bei Verwendung von Tabellen und auch Abbildungen beachten Sie bitte, dass diese immer Unter-/Überschriften enthalten (inklusive einer Nummer). Im Textfluss erklären/beschreiben Sie die Abbildung bzw. die Tabelle und nehmen Bezug über einen Verweis auf die Nummer.
\end{tcolorbox}


%Eizelne Kapitel
\newpage
\section{Benutzeroberfläche}\label{kap_benutzeroberflaeche}
\input{chapters/Benutzeroberfläche}

\newpage
\section{Serveranfragen}\label{kap_serveranfragen }
\subsection{GraphQL}
\paragraph{}
GraphQL ist eine Abfragesprache und Server-Laufzeitumgebung für APIs.
Ihre Aufgabe ist es, genau die Daten zu liefern, die anfordert werden, und nicht mehr.
\\
Mit GraphQL sind APIs schnell, flexibel und einfach für Entwickler.
\\ \\
Laut dem 2020 State of the API Report von Postman.com steht GrapQL an fünfter Stelle der spannendsten Technologien für 2021.
{Vgl. u.a. \cite{PM1}}

Im Hinblick auf die Art und Weise, wie Abfragen an den Server mithilfe von 
\\ GraphQL behandelt werden können, sind folgende Aspekte zu beachten.
\\
\begin{quote}
  \textbf{Vorteile}\\
  \begin{itemize}
    \item
          GraphQL-Aufrufe werden in einem einzigen Round Trip gehandhabt. Wir bekommen genau die Daten, die angefragt haben (kein Over-Fetching).

    \item
          Stark definierte Datentypen verringern das Risiko einer Fehlkommunikation zwischen Client und Server.

    \item
          GraphQL ist introspektiv. So können wir eine Liste der verfügbaren Datentypen anfordern. Dies ist ideal für automatisch erstellte Dokumente.

    \item
          Eine Anwendungs-API kann sich mit GraphQL weiterentwickeln, ohne dass bestehende Anfragen beeinträchtigt werden.
    \item
          GraphQL schreibt keine spezifische Anwendungsarchitektur vor. Es kann auf einer vorhandenen REST-API installiert und mit aktuellen API-Management-Tools verwendet werden.
    \item
          Als Alternative zu REST ermöglicht GraphQL Entwicklern die Erstellung von Abfragen zur Extraktion von Daten aus mehreren Quellen mit einer einzigen API-Abfrage.

  \end{itemize}

  \textbf{Nachteile}
  \begin{itemize}
    \item
          Für Entwickler, die sich bereits mit REST-APIs auskennen, bedeutet GraphQL weiteren Lernaufwand.
    \item
          Mit GraphQL verschiebt sich die Funktionalität von Datenabfragen zur Serverseite, was zusätzliche Komplexität für Serverentwickler bedeutet.

  \end{itemize}

  \footnote{Vgl. u.a. \cite{RH1}}
\end{quote}

\subsubsection{GraphQL Playground}
Mit GraphQL Playground haben wir die Möglichkeit, alle Abfragen und Mutationen zu testen. Wir erhalten Zugriff auf relevante Informationen wie verfügbare Felder und deren Datentyp. Diese Informationen werden aktualisiert, wenn der Servercode geändert wurde. Dadurch wurde eine aktuelle API-Dokumentation gewährleistet. Für unser Projekt war es sehr praktisch und hat die Kommunikation als Entwickler effizienter gemacht.
\\
%\begin{center}
%\includegraphics[scale=0.60]{GraphQL_Playground}\label{fig:GraphQL_Playground}
%\end{center}
%\textbf{Abbildung \autoref{fig:GraphQL_Playground}:}
Genau derselbe Code wird für die Abfrage von Benutzerinformationen später verwendet.
\newpage

\subsection{Implementierung der Serveranfragen mit GraphQL}
Alle Abfragen und Mutationen wurden in einem separaten Ordner gesammelt.
Damit soll eine saubere Struktur des Codes gewährleistet werden.
Diese wurden für die spätere Verwendung in den React-Komponenten exportiert.

Mithilfe der Hooks useQuery bzw. useMutation von Apollo Client wurden die Lese- und Schreibabfragen durchgeführt.
\\
\textbf{Apollo Client}\\
Warum haben wir uns für Apollo entschieden? Was ist ApolloClient?
TODO
%Coming...
\newpage

\subsubsection{Leseabfrage}
Nachdem eine Abfrage exportiert wurde, ist sie bereit, in einer React-Komponente \\
importiert und angewendet zu werden.

\begin{lstlisting}
import { GET_MY_INFO } from "./GraphQL/Queries"
import { useQuery } from "@apollo/client"

const { loading, error, data } = useQuery(
GET_MY_INFO,
ContextHeader(token),
)
//Code-Auszug in frontend/src/App.js

\end{lstlisting}
Die Konstante ContextHeader enthält das Token in der Struktur, die erforderlich ist, um die Abfrage nur dann stellen zu können, wenn der Benutzer dazu berechtigt ist.
\\
Sollte das Token einen undefinierten Wert, null oder ungültig enthalten, wird der Server ein Fehler zurückgegeben.
\\
Der useQuery Hook liefert ein Ergebnisobjekt, welches eine der folgenden Optionen zurückgibt.
\\\\
\textbf{loading:}\\
Ein boolescher Wert, der angibt, ob die Abfrage in Bearbeitung ist.
Wenn loading wahr ist, ist die Anfrage noch nicht abgeschlossen. Typischerweise kann diese Information verwendet werden, um einen Lade-Spinner anzuzeigen.
\\\\
\textbf{error:}\\
Ein Laufzeitfehler mit den Eigenschaften von GraphQL Errors und network Error.
Dieses enthält Informationen darüber, was bei der Abfrage fehlgeschlagen  ist.
\\\\
\textbf{date:}\\
Ein Objekt, das das Ergebnis der GraphQL-Abfrage enthält.
\\Es enthält die tatsächlichen Daten vom Server.
\\
\newpage

%TODO: AUCH MUTATIONS ZEIGEN
\subsubsection{Mutationen}
TODO

\newpage
\subsection{Axios}
Zusätzlich zu den GraphQL-Abfragen wurde eine Post-Anfrage mit Axios bereitgestellt.
Mit dieser war es möglich, Bilder auf die S3 Speicher von AWS hochzuladen.

\begin{quote}
  Axios ist ein Promise-basierter HTTP-Client für node.js und den browser. Auf der Server-Seite wir das modul http verwendet, während im Browser XMLHttpRequests (ajax) ausgeführt werden.
\end{quote}\footnote{Vgl. u.a. \cite{AX1}}

\begin{lstlisting}
    function disableBtn() {
        const uploadBtn = document.getElementById("uploadBtn")
        uploadBtn.disabled = true
        uploadBtn.style.background = "#000000"
      }
    
      function fileUploadHandler() {
        errored ? 
         disableBtn()
         /*
        Tritt ein Fehler auf, wird die Abfrage nicht gesendet und der Knopf deaktiviert 
        Ein Fehler tritt auf, wenn die Datei zu gross ist oder oder wenn sie ein ungueltiges  Format hat
        */
        :
        console.log("uploading pic...", file?.name)
        const fd = new FormData()
        /* avatar ist der Name der Datei, die hochgeladen wird
          file ist der zu sendende  Wert */
        fd.append("avatar", file)  
    
        axios
          .post(urlAvatar, fd, {
            headers: {
              "x-auth-token": TOKEN,
            },
          })
          .then((res) => {
            setState((state) => ({ ...state, avatar: res?.data?.location }))    
            //In location finden wir die URL fuer das gerade hochgeladende Bild 
          })
      }\end{lstlisting}

%Das HTML-Eingabefeld „input“ wurde folgendermaßen definiert.
%\begin{lstlisting}
 %   <input type="file" onChange={fileSelectedHandler} />
%\end{lstlisting}


Das Format der hochzuladenden Dateien wurde limitiert, damit nur zulässige Dateien an den Server gesendet werden.
\begin{lstlisting}
    const admittedImageFormats = ["png", "jpg", "jpeg"]
\end{lstlisting}

Die Größe der hochzuladenden Datei wurde um 1 MB abgegrenzt.
Durch die Eigenschaft „size“ der ausgewählten Datei konnten wir auf die Größe der Datei zugreifen.
\begin{lstlisting}
    const imageSize = e.target.files[0].size
\end{lstlisting}

\newpage
\section{Qualitätssicherung }\label{kap_QS}
\input{chapters/Qualitätssicherung}
\newpage

%Temporär, danach sollten wir das Package glossaries benutzen
\section{Glossar}\label{kap_glossar}
%\newglossaryentry{kiln}
%{
 % name=kiln,
  %description={German: Brennofen (m.);\\Français: fourneau (m.)},
  %plural=kilns
%}

%Make Glossary properly...
%\acrodef{VB}{Visula Basic}

\textbf{Cloud-Computing:}\\
...
\\\\
\textbf{Cloud-Dienste:}\\
...
\\\\
\textbf{On-Demand:}\\
...
\\\\
\textbf{On-Premise:}\\
...
\\\\
\textbf{Region:}\\
Die Region ist ein völlig unabhängiges und eigenständiges geografisches Gebiet. Jede Region hat mehrere, physisch getrennte und isolierte Standorte, die als Availability Zones bekannt sind. Beispiele für Regionen sind London, Dublin, Sydney, usw \footnote{\cite{AWS1}, Seite 42}.
\\\\
\textbf{Availability Zone:}\\
Eine Verfügbarkeitszone ist einfach ein Datenzentrum oder eine Sammlung von Datenzentren. Jede Verfügbarkeitszone in einer Region verfügt über eine separate Stromversorgung, Netzwerk und Konnektivität, um die Gefahr eines gleichzeitigen Ausfalls in beiden Zonen zu verringern \footnote{\cite{AWS1}, Seite 42}.
\\\\

\textbf{Instance family:}\\
Instanzfamilien sind eine Sammlung von EC2-Instanzen, die nach dem Verhältnis von Speicher, Netzwerkleistung, CPU-Größe und Speicherwerten zueinander gruppiert sind. Zum Beispiel bietet die m4-Familie von EC2 eine ausbalancierte Kombination von Rechen-, Speicher- und Netzwerkressourcen. \footnote{\cite{AWS1}, Seite 95}.
\\\\
Instagram-Story
\\\\
Tag
\\\\
Buckets
\\\\
PAYG %https://www.nimbix.net/glossary/pay-go
\\\\
Metadaten
\\\\
Startkonfiguration %https://docs.aws.amazon.com/autoscaling/ec2/userguide/create-launch-template.html
\\\\
Scale-In/Out
%\\\\
%Governance, Compliance ?
%\\\\
\newpage

\section{Zusammenfassung und Ausblick}\label{kap_zusammfAusbl}
(To-Do:)
\\Kapitelweise Kurzdarstellung der Inhalte (inklusive Referenzierung auf die \\Kapitelnummerierung) => Nach dem Motto: \textit{Was wurde wo beschrieben?}
\\Kurzdarstellung \textit{Problem – Lösungsweg – Ergebnisse}
\\Rückkopplung auf die Einleitung: Wurde die Zielstellung der Arbeit und die \\Fragestellung zufriedenstellend beantwortet?
\\Kritische Bewertung (sofern nicht bereits im Hauptteil geschehen)
\\Offene Probleme
\\Richtung der zukünftigen/möglichen Arbeiten
\\Erläuterung, warum welche Aspekte in der Arbeit nicht erläutert 

\subsection{Umweltbezogene Aspekte}
Esta tesis habla sobre monitoreo y optimizacion de recursos de manera financiera. Pero esas dos areas enfocadas a la economia, tiene impacto en el medio ambiente por las emisiones generadas por las granjas de servidores. 
Estadisticas dicen que en Europa/Alemania se generan x toneladas de CO2 provenientes de centros de computo. Por tanto al monitoreat y reducir cosotos, se estan evitando despilfarros y al final tambien emisiones de CO2. 
\\
\subsection{Test von den Werkzeugen und Maßnahmen}
Da es in dieser Arbeit zeitlich nicht gelungen ist, die Überwachungswerkzeuge und Optimierungsmaßnahmen umzusetzen, bleibt es noch sie in einer echten Umgebung zu testen. Es wäre möglich zu verifizieren, ob die hier genannten Maßnahmen zur vergleichbaren Einsparungen führen, wie die vom Cloud-Anbieter Amazon genannten.

Amazon bietet ein kostenloses Kontingent an, die jedoch für diese Tests nicht genug war. 
\\
\subsection{Bewusstsein in der gesamten Organisation}
Zusätzlich zu den bisher genannten Maßnahmen ist es wichtig, dass Verbraucher von Cloud-Diensten Bewusstsein für die Entstehung von Kosten entwikclen.[ODER sensibilisier werden?] Von dem Entwickler bis zum der IT-Manager, jeder sollte wissen, dass es so einfach ist, Cloud-Dienste mit ein paar Klicks zu beauftragen. Diese können in kurzer Zeit ungewünschte  Kosten verursachen oder sogar über Jahre hinweg wirtschaftliche Schäden verursachen. 
%Wer wird benachrichtigt. Welche Tools? https://docs.aws.amazon.com/de_de/AWSEC2/latest/UserGuide/monitoring_ec2.html
\\
\subsection{Die richtige Personen(Owneship verbreiten)}
Die technischen Maßnahmen zur Überwachung und Kostenreduzierung wurden dargelegt, aber jemand muss diese Analysen, Anpassungen und Entscheidungen durchführen. 
Deshalb ist es wichtig, bestimmte Personen zu berücksichtigen, die die Verantwortung für das Geschehen in den Cloud-Systemen übernehmen. Idealerweise Menschen, die sich für das Thema interessieren und über die notwendigen Kenntnisse verfügen, um die gesetzten Ziele zu erreichen. 
%Sie redete über DIESES t.ly/XJ24
\\
\subsection{5G is comming}
Mit 5G ist pronostiziert, dass mehr Daten[WIE VIELE AN WELCHEM JAHR?] automatisch von Maschinen produziert werden.
\\
\subsection{Langfristige Einsparungen sollten größer als Investitionen für Optimierung sein}
Kostenoptimierung UND -Überwachung SOLLEN DIE Einsparungen NICHT ÜBERSCHREITEN . 
TRUSTED ADVISOR NICHT FÜR JEDE FIRMA.
%https://content.aws.training/video/cmcfrm/de/x2/1.0.0/jwplayer.html?endpoint=https%3a%2f%2flrs.aws.training%2fTCAPI%2f&auth=Basic%20OjUzYmEwYTZmLTk0ZmMtNDAwZi1hODBlLWQ1YzA5NmNkOWY1MA%3d%3d&actor=%7b%22objectType%22%3a%22Agent%22%2c%22name%22%3a%5b%22zEjHPzGX10miDWp26Y_cLg2%22%5d%2c%22mbox%22%3a%5b%22mailto%3alms-user-zEjHPzGX10miDWp26Y_cLg2%40amazon.com%22%5d%7d&registration=2f22bc75-44b9-4175-a976-5ba4d7fe2902&activity_id=http%3a%2f%2fid.tincanapi.com%2factivity%2ftincan-prototypes%2fgolf-example&grouping=http%3a%2f%2fid.tincanapi.com%2factivity%2ftincan-prototypes%2fgolf-example&content_token=2554ffa1-5a1e-4737-9a0f-fc3abda1083e&content_endpoint=https%3a%2f%2flrs.aws.training%2fTCAPI%2fcontent%2f&externalRegistration=CompletionThresholdPercent%7c80!InstanceId%7c0!PackageId%7ccmcfrm_de_x2_1.0.0!RegistrationTimestampTicks%7c16324112989178350!SaveCompletion%7c1!TranscriptId%7cl4fipkeHAkKh-aGlnjYdug2!UserId%7czEjHPzGX10miDWp26Y_cLg2&externalConfiguration=&width=1366&height=728&left=0&top=0
%2:25
\begin{comment}
Von Buch "Gestaltung"
  Schluss (Fazit)
Den Abschluss der Arbeit bildet die Zusammenfassung der wesentlichen
Ergebnisse, die folgende drei Punkte beinhaltet:
Beantwortung der Forschungsfrage, die Sie in der Einleitung
aufgeworfen haben.
Sinnstiftung der Arbeit: Für welchen Zweck sollen die Ergebnisse
verwendet werden?
Gegebenenfalls auch persönliche Bemerkungen und Bewertungen oder
ein kurzer Ausblick.
\end{tcolorbox}

\end{comment}

%FAZIT ist (Von Schribe)
%Was ist? Hier sollten die neue Erkenntnisse der Arbeit dargestellt wrden

%TIPPS: 
%Vermeide "man" und "ich"
%Addressanten sind potenzielle Cloud-Nutzer/IT-Personal.
%Es wurde AWS ausgewählt, weil...
%Wirtschaftliche Betrachtung, weil das wichtig für Firmen ist

%STRUKTUR:
%Zusammenfassung, von was gemacht wurde?
%Beanwortung der FF 
%Mehrwer für die Praxis
%Limitationen: es konnte in dieser Arbeit nicht in der Praxis geprüft werden, ob die Maßnahmen ihre Verprechen anhalten
%Weitere Forschungen: es empfehlt sich diese Maßnahmen in echte Systeme einzusetzen

\newpage

\section{Literaturverzeichnis}\label{kap_literaturverzeichnis}
% INFO: Biblatex -Ausgabe des  
% Literaturverzeichnisses (Beispiele):   
% - \printbibliography => Ausgabe ALLER 
%   Einträge
% - \printbibliography[nottype=online]
%   => Ausgabe der Einträge, bis auf die
%      "Online"-Einträge
% - \printbibliography[type=online]     
%   => Ausgabe nur der "Online"-Einträge  
%\printbibliography


% Literaturverzeichnis
% INFO: Referenzieren auf das Literaturverzeichnis:
%
% Befehl: \cite{refmarke}
% 
% "refmarke" ist die Angabe in den geschweiften Klammern bei 
% \bibitem[]{refmarke}. 
\newpage
\thispagestyle{empty}
\section{Quellenverzeichnis}
\subsection{Literatur}
\renewcommand{\refname}{} % Literaturverzeichnis ohne Bezeichnung
% Literaturverzeichnis
\begin{thebibliography}{SW11} % 2. {...} => Hier die größte /breiteste Nummer (z.B. 99) oder Kurzbeleg angeben.
  \bibitem{SW11} Stickel-Wolf, Christine; Wolf, Joachim (2011): Wissenschaftliches Lernen und Lerntechniken. Erfolgreich studieren–-gewusst wie!. Wiesbaden: Gabler.
  % TODO cite correctly
  \bibitem{PG01} Seite 51 Zeile 5-6 \url{https://www.researchgate.net/profile/Ciprian-Octavian-Truica/publication/264416935_Asynchronous_Replication_in_Microsoft_SQL_Server_PostgreSQL_and_MySQL/links/53dbe6160cf216e4210c0375/Asynchronous-Replication-in-Microsoft-SQL-Server-PostgreSQL-and-MySQL.pdf}
  \bibitem{DB01} Scalability Databases \url{https://ieeexplore.ieee.org/abstract/document/7369245}
\end{thebibliography}

\subsection{Internetquellen}
\begin{thebibliography}{HR08} % 2. {...} => Hier die größte/breiteste Nummer (z.B. 99) oder Kurzbeleg angeben.
  \bibitem{BBoJ}Bertelsmeier, Birgit (o. J.): Tipps zum Schrei\-b\-en ei\-n\-er Ab\-sch\-luss\-ar\-beit. Fach\-hoch\-schu\-le Köln-Campus Gummersbach, Institut für Informatik. \url{http://lwibs01.gm.fh-koeln.de/blogs/bertelsmeier/files/2008/05/abschlussarbeitsbetreuung.pdf} (29.10.2013).
  \bibitem{HR08} Halfmann, Marion; Rühmann, Hans (2008): Merkblatt zur Anfertigung von Projekt-, Bachelor-, Master- und Diplomarbeiten der Fakultät 10. Fachhochschule Köln-Campus Gummersbach.\url{http://www.f10.fh-koeln.de/imperia/md/content/pdfs/studium/tipps/anleitungda270108.pdf} (29.10.2013).
  \bibitem{V01} Offizielle Vue-Website: Vergleich zwischen Vue, React und Angular. \url{https://vuejs.org/v2/guide/comparison.html#Preact-and-Other-React-Like-Libraries} (unbekannte Veröffentlichung).
  \bibitem{R01}Offizielle React-Website: React Hooks. \url{https://reactjs.org/docs/hooks-faq.html#which-versions-of-react-include-hooks}
  \bibitem{R02}Offizielle React-Website: Fortgeschrittene Anleitungen. \url{https://de.reactjs.org/docs/context.html}
  \bibitem{A01}Offizielle Website Apollo für React. \url{https://www.apollographql.com/docs/react/}
  \bibitem{SO01}StackOverFlow: Developer Survey 2021. \url{https://insights.stackoverflow.com/survey/2021#section-most-popular-technologies-web-frameworks}
  \bibitem{EE1}Elad Elrom: React and Libraries. \url{https://link.springer.com/content/pdf/10.1007%2F978-1-4842-6696-0.pdf}
  \bibitem{SS1}Stoyan Stefanov: Durchstarten mit React. \url{https://content-select.com/media/moz_viewer/5d5fc360-478c-4038-ac17-246bb0dd2d03/language:de}
  \bibitem{RH1}Red Hat: Was ist GraphQL? \url{https://www.redhat.com/de/topics/api/what-is-graphql}
  \bibitem{PM1}Postman: 2020 State of the API Report \url{https://www.postman.com/state-of-api/the-future-of-apis/#the-future-of-apis}
  \bibitem{AX1}Offizielle Website Axios \url{https://axios-http.com/}

  \bibitem{PG11}Offizielle Webseite PostgreSQL \url{https://www.postgresql.org/}
  \bibitem{PG12}PostgreSQL: Warum sollte man PostgreSQL verwenden? \url{https://www.postgresql.org/about/}
  \bibitem{PG13}Stackshare: Wer benutzt PostgreSQL? \url{https://stackshare.io/postgresql}

  \bibitem{MG1}Offizielle Webseite MongoDB \url{https://www.mongodb.com/de-de}
  \bibitem{MG2}Sharding \url{https://docs.mongodb.com/manual/sharding/}
  \bibitem{MG3}Replizierung \url{https://docs.mongodb.com/manual/replication/}
  \bibitem{MG2}NoSQL Erklärt \url{https://www.mongodb.com/de-de/nosql-explained}
  \bibitem{JSON1}Einführung in JSON \url{https://www.json.org/json-de.html}

  \bibitem{12FA1}The Twelve-Factor-App: X.Dev-Prod-Vergleichbarkeit \url{https://12factor.net/de/}


\end{thebibliography}
\newpage

\setcounter{section}{0} % Nummerierung der Gliederungsebene "section" auf 0 setzen
\renewcommand*\thesection{\Alph{section}} % Nummerierungsart für die Gliederungsebene "section" 
% auf Großbuchstaben setzen
\section{Anhang}\label{anhang}
\subsection{Aufwandsverteilung}\label{subsec_UabsAnhang}
Hier zeigen wir, wie die Aufgaben unter den Autoren des Projekts verteilt wurden.
TABELLE/Bild KOMMT...

\subsection{ANHAND X}\label{subsec_UabsAnhang}


\subsection{Verwendete Technologien und Werkzeuge}\label{subsec_UabsAnhang}
NodeJS
\\
React
\\
Bootstrap
\\
GitHub
\\
Cypress
\\
VSC
\\
Heroku
\\
LaTeX
\\
Cypress

\newpage

% Erklärung über die selbständige Abfassung der Arbeit  
\pagestyle{empty}
\section*{Erklärung über die selbständige\\Abfassung der Arbeit} % \section*{...}: das *-Symbol erlaubt, dass dieser
% Gliederungspunkt nicht ins Inhaltsverzeichnis aufgenommen wird
\addcontentsline{toc}{section}{Erklärung über die selbständige Abfassung der Arbeit}
Ich versichere, die von mir vorgelegte Arbeit selbständig verfasst zu haben.
Alle Stellen, die wörtlich oder sinngemäß aus veröffentlichten oder nicht veröffentlichten Arbeiten anderer entnommen sind,
habe ich als entnommen kenntlich gemacht.\\
Sämtliche Quellen und Hilfsmittel, die ich für die Arbeit benutzt habe, sind
angegeben. Die Arbeit hat mit gleichem Inhalt bzw. in wesentlichen Teilen noch keiner anderen Prüfungsbehörde vorgelegen.\\\\
\begin{tabular}{cp{7cm}}
                                    &             \\
                                    &             \\ \hline
  \small (Ort, Datum, Unterschrift) & \normalsize \\
\end{tabular}

%<MERKKASTEN> (für die eigene Verwendung bitte entfernen
\vspace{1cm}
\begin{tcolorbox}[title={Hinweise zur obigen \textit{Erklärung}}]
  \begin{itemize}
    \item Bitte verwenden Sie nur die Erklärung, die Ihnen Ihr \textbf{Prüfungsservice} vorgibt. Ansonsten könnte es passieren, dass Ihre Abschlussarbeit nicht angenommen wird. Fragen Sie im Zweifelsfalle bei Ihrem Prüfungsservice nach.
    \item Sie müssen \textbf{alle abzugebende Exemplare} Ihrer Abschlussarbeit unterzeichnen. Sonst wird die Abschlussarbeit nicht akzeptiert.
    \item Ein \textbf{Verstoß} gegen die unterzeichnete \textit{Erklärung} kann u.\,a. die Aberkennung Ihres akademischen Titels zur Folge haben.
  \end{itemize}
\end{tcolorbox}
%</MERKKASTEN>   

\newpage
% Unbeschriftetes Abschlussblatt (Leere Seite)
\thispagestyle{empty}
\input{leereSeite}

\end{document}

